\chapter{Conclusion, Recommendations and Future Work}\label{ch:conclusion}

This final chapter summarises the study, assesses the achievement of each objective, and turns the findings
into recommendations for management. It then states the contributions and limitations of the work, reflects on
the methodology, and proposes directions for future work.

\section{Summary of the study}
This project set out to design and develop an intelligent data analytics and machine learning framework that
transforms raw business data into actionable insights and predictive information. All seven objectives were
met. A dataset of \NRows transactions was validated, cleaned and stored in a dimensional data store (O1, O2).
Exploratory and inferential analysis revealed price-led growth, inelastic demand, ineffective discounts, strong
year-end seasonality and high geographical and portfolio concentration (O3). Seven regression and seven
classification algorithms, three clustering methods, statistical forecasting models and two anomaly detectors
were developed and compared under a strict time-ordered protocol (O4). Each of the nine algorithms named in the
proposal was also studied on its own, which explained its behaviour and independently confirmed the results.
The results are delivered through a nine-page interactive dashboard and a project website that are publicly
deployed free of charge, with continuous deployment from GitHub (O5), and through automatically generated
insights and recommendations (O6). The framework's predictive performance, usefulness, efficiency and reproducibility were
evaluated (O7).

The central conclusion is that the greatest value comes from \emph{integration} and \emph{discipline}, not from
the most complex algorithm. Clean and validated data, leakage-safe features, honest benchmarks and clear
communication turned simple and well-understood models into reliable decision support. The best
product-level forecasts cut the error of ``same as last year'' planning by \GainVsSNaive\%. Classical and
machine-learning methods worked best together. The organisation can thus move from describing the past to
anticipating the future, answering \emph{what happened}, \emph{why} and \emph{what next} within one system.

\section{Achievement of the objectives}
\Cref{tab:objectives} maps each objective of \cref{ch:intro} to the work that fulfilled it and to the evidence
in this report.

\begin{table}[htbp]\centering\small
\caption{Achievement of the research objectives}\label{tab:objectives}
\begin{tabularx}{\linewidth}{@{}lXll@{}}
\toprule
\textbf{Obj.} & \textbf{How it was achieved} & \textbf{Evidence} & \textbf{Status} \\ \midrule
O1 & \NRows transactions profiled, validated and documented in a data dictionary & \cref{tab:dqsummary,tab:dictionary} & Achieved \\
O2 & Defensive, tested preprocessing pipeline with a cleaning log & \cref{tab:tests-unit} & Achieved \\
O3 & Trends, seasonality, concentration, hypothesis tests and elasticity & \cref{ch:eda} & Achieved \\
O4 & Seven regressors, seven classifiers, clustering, anomaly detection and statistical benchmarks; nine-algorithm study & \cref{ch:models,ch:algorithms} & Achieved \\
O5 & Nine-page dashboard and website, publicly deployed & \cref{ch:dashboard,tab:links} & Achieved \\
O6 & Automatically generated insights with actions & \cref{sec:insights} & Achieved \\
O7 & Evaluation of accuracy, usefulness, efficiency and reproducibility & \cref{tab:syseval} & Largely achieved\textsuperscript{a} \\ \bottomrule
\end{tabularx}
\par\smallskip\footnotesize \textsuperscript{a}A formal usability study with managers remains to be done (\cref{sec:limitations}).
\end{table}

The research questions were all answered in \cref{ch:discussion}, and the hypothesis decisions are summarised
in \cref{tab:hypdecision}.

\section{Recommendations for management}
\begin{itemize}
  \item \textbf{Pricing and promotions.} Demand is inelastic and discounts do not raise order size. The discount
        budget of about \BDT\DiscountCost million over the period should be reviewed. Controlled experiments
        (A/B promotions) should establish whether targeted promotions pay off.
  \item \textbf{Volume growth.} Because growth relies on price, volume should become a separately tracked KPI,
        with growth initiatives in under-penetrated regions and in E-commerce.
  \item \textbf{Seasonal planning.} Production and distributor stock should be built from October for the
        \PeakMonth peak. February is the natural window for maintenance and promotions.
  \item \textbf{Portfolio.} The A-class products deserve the highest service levels. C-class products are
        candidates for review.
  \item \textbf{Logistics.} Large orders are concentrated in the distributor channel and in particular regions.
        Vehicle and stock allocation can be planned at that level.
  \item \textbf{Control.} A daily anomaly screen improves data quality and detects operational problems early.
\end{itemize}

\section{Contributions of the study}
\paragraph{Contribution to practice.} The study delivers a working, publicly deployed decision-support system and
a set of evidence-based recommendations for a real product portfolio. Because every component is open source
and configured in one place, the framework can be reused for other companies and domains at no licence cost.

\paragraph{Contribution to knowledge.} The study provides a documented comparison of classical forecasting,
seven machine-learning regressors and seven classifiers under one strict, time-ordered protocol on a
real-scale consumer-goods dataset from Bangladesh. It confirms, in a new context, that simple and statistical
models are hard to beat when validation is honest \citep{hyndman2021}, and it adds context-specific evidence
on price inelasticity, discount ineffectiveness and channel-driven order size.

\paragraph{Methodological contribution.} The combination of (i) generated reporting, in which no number is typed
by hand, (ii) an independent second implementation of the models, and (iii) automated tests for data leakage
is offered as a practical recipe for reproducible applied analytics, in line with the concerns raised by
\citet{sculley2015} and \citet{kaufman2012}.

\section{Limitations}\label{sec:limitations}
\begin{itemize}
  \item \textbf{Data provenance and realism.} The dataset is very clean: it has no missing values, its
        accounting identities hold exactly, and order size is almost independent of customer type, discount and
        product. Real operational data are usually messier. The pipeline was therefore tested on deliberately
        damaged copies, but conclusions about customer behaviour should be validated on further company data
        before they are acted on. The exact source and collection process of the dataset should be documented
        alongside this report.
  \item \textbf{Missing drivers.} No data were available on marketing spend, competitor prices, stock-outs,
        weather or holidays. Including them would probably improve both forecasting and classification.
  \item \textbf{Observational elasticity.} The elasticity estimate rests on list-price variation over time, not
        on randomised experiments. Prices changed by small amounts, so elasticity outside the observed range is
        unknown.
  \item \textbf{Future conditions.} All forecasts assume that past patterns continue. Structural breaks, such as
        a new competitor or an economic shock, would reduce accuracy.
  \item \textbf{Scope of evaluation.} The dashboard has not yet been evaluated with real managers. Deep-learning
        models were not explored, because the panel (52 series of 81 months) is too small to justify them.
\end{itemize}

\section{Reflection on the methodology}
Strict time-ordered evaluation with benchmarks was the most important methodological decision. Random
cross-validation would have produced flattering but misleading accuracy figures. Generating every reported
number from code turned out to be valuable as well: analyses were re-run many times as the work progressed, and
the report always remained consistent. Running a second, independent modelling study (\cref{ch:algorithms})
proved a cheap and effective way to test whether the conclusions depended on one implementation. Both practices
are recommended for future projects of this kind.
\section{Future work}
\begin{itemize}
  \item \textbf{Richer data:} add promotions, marketing spend, competitor prices, inventory and stock-outs,
        public holidays (Eid, Pohela Boishakh) and weather, and integrate them through APIs.
  \item \textbf{Real-time analytics and alerts:} schedule the pipeline daily in the cloud and send anomaly and
        forecast-deviation alerts by e-mail or messaging.
  \item \textbf{MLOps:} version data and models, monitor forecast accuracy over time and retrain automatically
        when performance drifts.
  \item \textbf{Prescriptive analytics:} optimise production quantities and distributor allocations under
        capacity constraints, using the forecast intervals.
  \item \textbf{Causal evaluation of promotions:} run controlled promotion experiments to measure the true
        effect of discounts.
  \item \textbf{Advanced models:} hierarchical reconciliation across product, region and channel, and global
        deep-learning forecasters once longer and more granular histories are available.
  \item \textbf{Generative-AI assistant:} let managers ask questions in natural language (for example,
        ``why did Rangpur sales fall in March?'') and receive answers grounded in the framework's data and metrics.
  \item \textbf{Usability study:} evaluate the dashboard with managers using task-based tests and the System
        Usability Scale \citep{brooke1996}.
\end{itemize}

\section{Concluding remarks}
The project began with a familiar situation: plenty of data, spreadsheet reports that describe the past, and
decisions made largely on experience. It ends with a reproducible framework that validates the data, explains
what drives sales, predicts demand with measured accuracy and delivers the results to managers through a
public dashboard with recommended actions. The most important lesson is that dependable decision support comes
less from sophisticated algorithms than from careful data preparation, honest evaluation and clear
communication. These principles, and the framework that embodies them, can be carried forward to future
projects in any organisation that wants to move from describing its past to anticipating its future.
